\begin{abstract}
Large language models (LLMs) are increasingly used as annotators in computational social science, yet annotation errors propagate into downstream causal estimates in ways that binary measurement error theory cannot predict. We present the first systematic framework characterizing how $K$-class confusion matrix structure interacts with DAG topology to determine bias direction and magnitude. We derive bias expressions for seven canonical DAG topologies under multi-class non-differential misclassification (closed-form for confounding, exposure, and mediation; numerically validated for all seven) and establish a structural sign-flip guarantee (with formal proofs for three topologies and Monte Carlo validation for all seven): five topologies (confounding, mediation, collider, M-bias, front-door) never produce sign reversal regardless of the confusion matrix, while exposure and instrumental variable topologies can. In the exposure DAG, 83.4\% of $K{=}3$ confusion matrices cause sign reversal, a failure that persists at 13.7\% violation rate (2.1\% sign reversal, 11.7\% amplification) even with classifier accuracy above 70\%. We introduce Annotation Sensitivity Values (ASVs), a sensitivity analysis framework with three complementary measures that reveal a topology-dependent vulnerability spectrum: IV estimates are maximally vulnerable (52.6\% sign reversal; sign-flip ASV $= 1.1$--$6.7$) while M-bias adjustment is immune (ASV $= \infty$). We evaluate five correction methods across all topologies and prove a Frisch--Waugh--Lovell invariance result: matrix-based corrections achieve 0\% bias reduction on the target coefficient for four topologies, while PPI++ achieves 93.8\% reduction but requires gold labels. Binarization from multi-class to binary introduces additional bias varying 147$\times$ across topologies. Neither classification $F_1$ nor correction method choice alone determines downstream safety; topology-aware analysis is necessary for valid causal inference with LLM-generated annotations.
\end{abstract}
